
\chapter{Жизнь транзакции: от касания карты до денег на счёте}\label{ch:tx-lifecycle}
Покупатель прикладывает карту в~кофейне, терминал показывает
\enquote{одобрено} -- но деньги продавец увидит на счёте позже,
иногда через несколько суток. Карточная транзакция -- многофазная
операция, и~на каждом шаге меняется соотношение факта и
обещания.

\section{Авторизация}\label{sec:tx-authorization}

Авторизация -- единственная фаза карточного платежа, которая идёт
в~реальном времени и~видна покупателю. Терминал показывает
\enquote{одобрено} или \enquote{отказано}, и~на этом контакт покупателя
с~платёжной системой заканчивается. Клиринг, расчёт, возможный диспут
случаются за кулисами, иногда через дни после покупки. Но именно
в~эти несколько секунд фиксируются все данные, по которым строится
вся дальнейшая обработка.

\subsection*{Путь авторизационного запроса}

Когда держатель прикладывает карту к~терминалу, вставляет её
или проводит магнитной полосой, терминал считывает данные
карты и формирует авторизационный запрос -- структурированное
сообщение в~формате ISO~8583; подробно о~формате см.\
в~главе~\ref{ch:iso8583}. Внутри такого сообщения собраны как
минимум PAN в~DE~2 и~данные Track~2 в~DE~35 или их чиповый
эквивалент в~DE~55, сумма и~валюта операции в~DE~4 и DE~49,
идентификаторы терминала и ТСП в~DE~41 и DE~42, категория
ТСП MCC в~DE~18, а~также тип операции, который задают MTI
и processing code в~DE~3.

Сообщение уходит эквайеру -- банку или процессору, обслуживающему
этот терминал. Эквайер выполняет первичные проверки -- валидность
терминала, соответствие MCC, базовые антифрод-правила на своей
стороне -- и маршрутизирует запрос в~карточную сеть~\cite{visa-core-rules}.
Сеть по IIN/BIN определяет эмитента или его процессор
и~передаёт запрос ему.

Весь путь -- терминал, эквайер, сеть, эмитент и~обратно --
должен уложиться в~жёсткие временны\'{е} рамки. У Visa публичные
правила задают верхние пределы ожидания для POS-авторизаций
порядка нескольких секунд, с~более коротким пределом в~Европе.
Таймауты для снятия наличных в~банкомате исторически задавались
отдельно, но в~редакции Visa Core Rules от 18~апреля 2026~года
требование к ATM-таймауту снято и оставлено на усмотрение
участников~\cite{visa-core-rules}. Если ответ не пришёл вовремя,
за эмитента может ответить сама Visa через механизм Stand-In
Processing (STIP). STIP одобряет только операции с~низким профилем
риска -- сумма в~пределах нормы для карты и~ТСП, знакомая география,
отсутствие сигналов фрода. Это консервативнее решения эмитента,
поскольку сеть предпочитает отклонить неясную операцию, а~не принимать риск
на себя. STIP не заменяет эмитента -- он покрывает только безопасный
хвост распределения транзакций. У Mastercard пределы ожидания
и показатели доступности задаются в~спецификациях сообщений
и контролируются через нормативы доступности авторизации~\cite{visa-core-rules,mc-tpr}.

\subsection*{Решение эмитента}

Получив авторизационный запрос, эмитент или его процессор
решает одобрить или отклонить операцию по нескольким группам
проверок, которые выполняются за доли секунды.

\begin{enumerate}
  \item \textbf{Проверка карты.} Существует ли карта с~таким PAN?
    Не заблокирована ли она? Не истёк ли срок действия?
    Совпадает ли CVV или криптограмма?
  \item \textbf{Проверка баланса.} Достаточно ли средств на счёте
    дебетовой карты или доступного кредитного лимита
    у кредитной?
  \item \textbf{Лимиты и ограничения.} Не превышен ли дневной лимит
    на операции? Не заблокирована ли карта для этого типа
    операций -- скажем, запрет на интернет-покупки? Не ограничена ли
    география использования?
  \item \textbf{Антифрод.} Соответствует ли операция типичному
    поведению держателя? Нет ли признаков компрометации карты?
    Не входит ли ТСП в~списки повышенного риска?
  \item \textbf{Аутентификация держателя.} Для чиповых транзакций
    верификация криптограммы ARQC; для CNP -- результат
    3-D Secure, если он был; для магнитной полосы -- CVV.
\end{enumerate}

По результатам этих проверок эмитент формирует ответ с~кодом
ответа в~DE~39. Значение \texttt{00} означает, что операция одобрена;
\texttt{05} -- общий отказ без детализации; \texttt{51} -- недостаточно
средств; \texttt{14} -- неверный номер карты, и~так далее. Полный
список кодов зависит от ISO~8583-профиля и~схемных документов;
базовая модель стандартизована, но сети дополняют её собственными
значениями и уточнениями~\cite{iso-8583-2023,visa-tadg,mc-rules}.

Если транзакция одобрена, эмитент присваивает ей код авторизации
в~DE~38 -- шестисимвольный идентификатор,
привязывающий одобрение к~конкретному запросу. На этапе клиринга
именно по нему сопоставят авторизацию с~финансовой записью.
Одновременно эмитент блокирует запрошенную сумму на счёте
держателя -- холд (\textit{hold}). Деньги ещё не списаны, но уже недоступны для других
операций.

\subsection*{Срок жизни холда}

Холд не бессрочен, но и~не живёт по единому универсальному таймеру.
Точнее говорить об окне между авторизацией и тем моментом, когда
транзакцию нужно либо представить в~клиринг, либо снять
отменой -- \textit{reversal}. У Visa для большинства транзакций с~физическим
предъявлением карты это окно составляет 5 календарных дней;
для операций без предъявления -- 10 дней; для оценочной авторизации
в~гостиницах, круизах и~аренде автомобилей -- до 30 дней.
У Mastercard клиринговое окно зависит от продукта и~региона
и~находится в~диапазоне порядка 7--30 дней
(пункт~3.15.1 TPR).
Эти окна и~задают, как долго эмитент держит неиспользованный
резерв; точный момент снятия холда зависит ещё и~от внутренних
правил эмитента~\cite{visa-core-rules,mc-tpr}.

Числа 5, 7 и 30 дней отражают минимальный операционный цикл -- время
доставки или оказания услуги плюс время передачи клирингового файла
от эквайера в~сеть плюс буфер на задержки между участниками.
Длинные окна для гостиниц и~аренды автомобилей в 30 дней
соответствуют реальной дельте между \textit{check-in}-авторизацией и финальной
суммой при выезде. Сроки сложились исторически и держатся консенсусом
схем; их изменение требует согласования всей расчётной цепочки.

\practicenote{%
  Если клиент жалуется, что \enquote{деньги списали, но покупка
  не прошла}, чаще всего речь идёт о~холде. Авторизация была
  одобрена и~средства заблокированы, но по какой-то
  причине продавец не завершил операцию -- скажем, отменил заказ
  на своей стороне, не отправив отмену. Деньги вернутся
  автоматически по истечении холда, но клиенту это не очевидно,
  и~он видит в~мобильном банке уменьшенный доступный баланс.
}

\practicenote{%
  Частичный отказ авторизационного хоста эмитента создаёт
  рассинхронизацию холда и реального состояния счёта. Если запрос
  был одобрен до отказа, но обновление таблицы холдов не сохранилось,
  сумма остаётся заблокированной на неопределённый срок, хотя
  клиент видит \enquote{отказ} в~мобильном банке. Противоположный сценарий --
  холд снят слишком рано при откате незафиксированной транзакции,
  и~клиент тратит средства, которые ещё \texttt{pending}. Для этого класса
  расхождений операционные команды ведут отдельный
  журнал записей холдов с~меткой времени авторизации и признаком сопоставления с~клирингом.

}

\section{Клиринг}\label{sec:tx-clearing}

Авторизация -- обещание заплатить; клиринг -- предъявление счёта,
формальное требование эквайера к~эмитенту это обещание исполнить.
Клиринг идёт не в~реальном времени, а~пакетами (batches), и~его
механика принципиально иная. Нет мгновенного ответа, нет покупателя,
ждущего результат, нет жёстких таймаутов в~секундах.

\subsection*{Формирование клирингового файла}

В~конце операционного дня, а~в зависимости от настроек и
несколько раз в~день, эквайер формирует клиринговый файл (\textit{clearing file}) --
пакет записей обо всех транзакциях, прошедших авторизацию за период.
Каждая запись содержит реквизиты транзакции -- PAN, сумму, которая
может отличаться от авторизованной, валюту, дату и~время,
идентификатор ТСП, MCC, код авторизации и другие данные,
нужные для расчёта~\cite{mc-rules}. Отправку такого файла с~итоговыми
суммами называют финансовым предъявлением (\textit{financial presentment})
или захватом (\textit{capture}); в~этот момент эквайер фиксирует финальную
сумму обязательства~\cite{visa-tadg,mc-tpr}.

Формат клирингового файла зависит от карточной сети. Visa использует
BASE II, в~котором каждый тип записи обозначается Transaction Code (TC);
Mastercard -- формат IPM (Integrated Product Messages);
НСПК -- собственный формат для внутрироссийских операций. Логика
у~всех одна -- набор записей с~реквизитами транзакций, по которым сеть
должна провести расчёт~\cite{visa-core-rules,mc-rules}.

\subsection*{Сопоставление с~авторизацией}

Получив клиринговый файл, сеть сопоставляет каждую запись
с~соответствующей авторизацией по комбинации PAN, суммы, кода
авторизации и даты. Сопоставление не всегда тривиально, потому что сумма
в~клиринге может отличаться от авторизованной. Типичный пример --
ресторан, где авторизация проходит на сумму заказа, а в~клиринг попадает
сумма с~чаевыми, которые клиент добавил позже. Правила схем допускают
такое расхождение только для определённых сценариев -- чаевых,
категорий travel\&entertainment, -- а~точные пределы зависят от MCC,
программы и текущей редакции правил~\cite{visa-core-rules,mc-rules}.

Если предъявление в~клиринг приходит позже применимого окна
или данные в~нём не бьются с~исходной авторизацией, сеть всё равно
может доставить это предъявление эмитенту. Но для эквайера это
зона позднего предъявления (\textit{late presentment}), где повышается риск диспута, и~эквайер теряет
часть защитного временного окна по спору~\cite{visa-core-rules,mc-tpr,mc-chargeback-guide}.

Острее сценарий, когда клиринговый файл теряется или не принимается
сетью из-за ошибки формата. Проблему обнаруживает эквайер, когда сеть
возвращает отчёт с~перечнем отклонённых записей -- \textit{rejected
transactions}. Эквайер чинит записи и~повторно отправляет файл
в~следующем батч-окне. Формально правила сетей ограничивают срок
повторной подачи. У Visa \textit{financial presentment} должен попасть в~сеть
в~пределах авторизационного окна -- 5--30 дней в~зависимости от MCC, --
иначе авторизация аннулируется~\cite{visa-core-rules}. У Mastercard
клиринг должен прийти в~пределах 7--30 календарных дней~\cite{mc-tpr}.
Если эквайер не уложился, эмитент автоматически снимает холд,
и~деньги возвращаются в~доступный баланс держателя. Для продавца
это значит, что операция прошла в~терминале, товар отдан, но через
карточную сеть деньги уже не получить, и взыскивать придётся
иначе. Поэтому статус каждого клирингового батча мониторится
эквайером отдельно.

Если же авторизации не существовало вовсе, эмитент всё равно
проводит клиринговую запись по счёту держателя без предварительного
hold'а. Деньги списываются в~момент поступления записи, а~не при
авторизации. Риск переходит к~эквайеру и~продавцу; именно этот
сценарий описывает force post в~разделе~\ref{sec:tx-edge-cases}~\cite{mc-tpr}.

\subsection*{Расчёт interchange и комиссий}

На этапе клиринга сеть рассчитывает interchange для каждой
транзакции -- ту самую межбанковскую комиссию, которую эквайер
платит эмитенту; см.\ главу~\ref{ch:ecosystem}. Ставка interchange
зависит от параметров, зафиксированных в~клиринговой записи, --
типа карты, MCC, способа приёма, географии операции и ряда
дополнительных признаков. Тип карты бывает дебетовый, кредитный
или корпоративный; способ приёма -- чип, магнитная полоса
или CNP; география -- внутренняя или трансграничная.

Параллельно сеть начисляет собственные комиссии -- плату за обработку,
за трансграничную маршрутизацию, за использование специальных
сервисов. В~схемных документах эту составляющую называют scheme fees
и assessments. По итогам обработки всех клиринговых записей за период
сеть формирует чистые расчётные позиции (net settlement positions)
для каждого банка-участника -- итоговые суммы, которые один банк
должен другому с~учётом взаимозачёта встречных потоков.

\section{Межбанковский расчёт (settlement)}\label{sec:tx-settlement}

Расчёт завершает жизненный цикл транзакции -- именно здесь деньги
действительно двигаются между банками. До этого момента средства
не покидали счёт держателя; авторизация лишь заблокировала сумму,
клиринг сформировал требование. На этапе расчёта требование
исполняется в~расчётной инфраструктуре платёжной
системы~\cite{fz-161,visa-core-rules,mc-tpr}.

\subsection*{Нетто-расчёт}

Карточные сети используют нетто-расчёт. Вместо отдельной
межбанковской операции на каждую транзакцию сеть суммирует все
встречные обязательства между парами банков за расчётный период
и переводит только разницу. Если за день банк А~как эквайер должен
банку Б как эмитенту миллион рублей interchange, а~банк Б как
эквайер должен банку А~как эмитенту восемьсот тысяч, фактически
переводятся только двести тысяч от А~к~Б. Этот механизм
многократно сокращает количество и объём межбанковских проводок.

Сеть проводит расчёт через назначенный расчётный банк или иной
расчётный сервис, предусмотренный её правилами. У Visa локальные
операции, если в~стране доступен соответствующий сервис, проходят
через National Net Settlement Service (NNSS), остальные --
через International Settlement Service; правила отдельно определяют
расчётный банк и расчётный счёт участника. У Mastercard итоговый
документ дня -- ежедневное уведомление по нетто-расчётам, и~именно оно
считается окончательной записью о~движении средств между
участниками~\cite{visa-core-rules,mc-tpr}.

\subsection*{Когда продавец получает деньги}

Для продавца расчёт начинается в~момент, когда деньги за проданный
товар или оказанную услугу поступают на его расчётный счёт. Между
расчётом в~карточной сети и зачислением на счёт продавца есть ещё
один шаг. Эквайер, получив средства от сети, перечисляет их продавцу
по условиям договора эквайринга -- ежедневно, через день или
по иному графику, за вычетом MSC.

Универсального срока здесь нет: после расчёта в~сети эквайер
перечисляет деньги продавцу по графику договора, и~на практике
это следующий рабочий день или иной срок по SLA конкретного
провайдера. Для продавца с~большим оборотом или узкой маржой каждый
дополнительный день -- замороженные оборотные средства, поэтому
скорость расчёта часто становится предметом переговоров с~эквайером.

\begin{figure}[tbp]
  \centering
  \includefiguresvg[width=\linewidth]{assets/figures/ch04-transaction-lifecycle/tx-lifecycle-timeline}
  \caption{Жизненный цикл карточной транзакции: от касания карты до
  зачисления продавцу.}\label{fig:tx-lifecycle-timeline}
\end{figure}

\section{Одностадийная и двухстадийная схема}\label{sec:tx-dual-single}

Описанная выше схема -- двухстадийная, dual message system (DMS):
авторизация и клиринг идут разными сообщениями в~разное время.
Сначала терминал спрашивает \enquote{можно?}, потом эквайер
предъявляет финансовую запись. Это стандарт для большинства
карточных покупок~\cite{visa-core-rules,mc-tpr}.

Существует и другая архитектура -- одностадийная, single message
system (SMS), -- в~которой авторизация и финансовая запись совмещены в~одном
сообщении. Такой запрос одновременно спрашивает разрешение
\emph{и} фиксирует финансовое обязательство. Ответ эмитента означает
сразу \enquote{одобрено} и \enquote{деньги можно считать списанными},
а отдельного клирингового шага не требуется.

\subsection*{Где используется каждая схема}

Двухстадийная схема доминирует в~обычных карточных покупках --
в~магазинах, ресторанах, во многих интернет-операциях, -- потому
что допускает расхождение между авторизованной и финальной суммой.
Это важно там, где возможны чаевые, довзвешивание товара или
частичная отмена. Кроме того, эквайер группирует транзакции в~пакеты
и обрабатывает их не по одной, а~сериями.

Одностадийная схема -- классическая модель для банкоматов (ATM).
Когда держатель снимает наличные, разделять авторизацию и финансовую
запись бессмысленно: деньги уже выданы физически, и~задним числом
операцию не отменить. По той же логике одностадийная схема используется
в~некоторых дебетовых сетях, где транзакция всегда проходит на точную
сумму и последующей корректировки не предполагает~\cite{visa-core-rules,mc-tpr}.

На уровне ISO~8583 различие выражается в MTI -- Message Type
Indicator. Авторизационный запрос в~двухстадийной схеме имеет MTI~0100,
а финансовый запрос в~одностадийной схеме -- MTI~0200;
соответствующие ответы -- MTI~0110 и MTI~0210. Структура MTI подробно
разбирается в~главе~\ref{ch:iso8583}.

\importantnote{%
  Выбор между двухстадийной и одностадийной схемой принимает не продавец
  или разработчик: это свойство канала и~правил карточной сети.
  POS-терминал, работающий по протоколу Visa или Mastercard,
  отправляет двухстадийный поток; банкомат, работающий в~сети ATM,
  отправляет одностадийный. Попытка отправить финансовое сообщение MTI~0200
  туда, где ожидается авторизационное MTI~0100, приведёт
  к~отказу сети.

}
\subsection*{Гибридные варианты}

На практике различие между двухстадийной и одностадийной схемой
размывается. Visa, Mastercard и процессоры поверх них поддерживают
операционные гибриды, в~которых одобренная авторизация при
определённых условиях автоматически доводится до финансовой записи
или обслуживается общей инфраструктурой авторизации и клиринга.
Гибриды упрощают жизнь эквайерам, но добавляют сложности
процессинговым системам, которым приходится держать оба сценария
в~одном коде.

\section{Три способа отменить платёж}\label{sec:tx-cancellation}

Отмена карточного платежа -- область, где разрыв между авторизацией
и расчётом создаёт максимальную сложность. Банковский перевод
либо прошёл, либо нет; карточная транзакция в~каждый момент времени
находится в~одной из нескольких фаз, и~способ отмены жёстко
определяется фазой.

\subsection*{Сообщение отмены (reversal)}

Отмена -- reversal -- аннулирует авторизацию \emph{до того}, как
транзакция ушла в~клиринг. В~терминах ISO~8583 это отдельный тип
сообщения. MTI~0400, Reversal Request, требует ответа от эмитента;
MTI~0420, Reversal Advice, односторонне уведомляет о~состоявшейся
отмене. Эквайер отправляет одно из них в~сеть с~просьбой
аннулировать ранее одобренную авторизацию.

Отмена возникает в~нескольких типовых сценариях. Покупатель передумал
у~кассы; терминал не получил ответ на авторизацию по таймауту
и~на всякий случай отправляет отмену; кассир ошибся с~суммой
и~отменяет операцию. При успешной отмене эмитент снимает холд
со счёта держателя, и~средства возвращаются в~доступный баланс
обычно мгновенно или за несколько минут. Правила Mastercard
формулируют обязанность жёстко -- эквайер должен обеспечить отправку
Reversal Request/0400 в~течение 24~часов после отмены ранее
авторизованной операции или её завершения на меньшую сумму;
у Visa отмена также обязательна, если одобрение пришло уже после
таймаута~\cite{visa-core-rules,visa-tadg,mc-tpr}.

Отмена -- самый прямой способ аннулирования. Работает в~реальном
времени, не порождает финансовой записи, не требует возврата денег,
потому что деньги ещё не перемещались. Но окно её применимости
закрывается с~уходом транзакции в~клиринг -- попавшую в~клиринговый
файл авторизацию reversal уже не отменит~\cite{mc-rules}.

\subsection*{Аннулирование до расчёта (void)}

\textit{Void}, или аннулирование, -- не отдельный тип сообщения карточной сети,
а операционный термин эквайринга и кассового ПО для отмены операции
до окончательного клиринга. Visa описывает сценарий прямо. В
\textit{dual-message}-среде \textit{cancellation} должна происходить до того, как
транзакция уйдёт в~клиринг на сторону эквайера; если отмена
случилась уже после выхода на онлайн-авторизацию, устройство или продавец
обязаны сформировать отмену, а~саму транзакцию убрать из пакета
клиринга или пометить как \textit{void}~\cite{visa-tadg}.

Для держателя карты \textit{void} и~отмена выглядят одинаково -- холд снимается,
деньги возвращаются на счёт. На технической стороне \textit{void} обрабатывается
медленнее, потому что он привязан к~циклу клиринга, и~если файл уже отправлен,
корректировка попадёт только в~следующий батч.

\practicenote{%
  В API большинства PSP и платёжных шлюзов вы встретите метод
  \texttt{cancel} или \texttt{void}, который на самом деле
  абстрагирует оба механизма. Если транзакция ещё не ушла
  в~клиринг, PSP отправляет отмену; если ушла, но расчёт
  не произошёл, применяется \textit{void}. Для разработчика, интегрирующего платежи,
  это удобная абстракция, но внутри неё скрываются разные
  процессы с~разной скоростью возврата средств клиенту.

}
\subsection*{Возврат после расчёта (refund)}

Возврат -- единственный способ вернуть деньги после того, как
расчёт завершён и~средства фактически перемещены между банками.
В~отличие от отмены и \textit{void}, которые аннулируют исходную транзакцию,
возврат -- \emph{новая} транзакция в~обратном направлении: эквайер
инициирует перевод от продавца к~держателю через ту же карточную
сеть.

Возврат проходит собственный цикл авторизации, клиринга и~расчёта,
только в~обратную сторону. В Visa TADG возврат описан как
онлайн-авторизационное сообщение с~последующим клирингом; Mastercard
прямо допускает возврат как операцию в \textit{dual-message}- или
\textit{single-message}-режиме, а~для возвратов в \textit{dual-message}-режиме
требует онлайн-авторизацию~\cite{visa-tadg,mc-tpr}.

Для держателя карты возврат -- самый медленный вариант: средства
поступают на счёт через несколько рабочих дней, потому что возврат
проходит полный цикл клиринга и~расчёта. Для продавца возврат
обычно и~самый дорогой операционно: часть расходов исходной покупки
автоматически не компенсируется, а~эквайер нередко берёт отдельную
комиссию за саму операцию возврата.

\begin{figure}[tbp]
  \centering
  \includefiguresvg[width=\linewidth]{assets/figures/ch04-transaction-lifecycle/tx-cancellation-timeline}
  \caption{Окна применимости отмены, \textit{void} и~возврата на временной
  шкале транзакции.}\label{fig:tx-cancellation-timeline}
\end{figure}

\importantnote{%
  Распространённая ошибка при проектировании платёжной интеграции --
  использовать возврат там, где достаточно отмены или \textit{void}.
  Если клиент отменяет заказ через минуту после оплаты, а~система
  продавца отправляет возврат вместо отмены, клиент будет ждать
  возврата несколько дней вместо нескольких минут, продавец потеряет
  interchange, а~эквайер зафиксирует лишнюю транзакцию. Правило
  простое -- отменяй как можно раньше; отмена лучше \textit{void},
  \textit{void} лучше возврата.

}
\section{Граничные случаи}\label{sec:tx-edge-cases}

Схема \enquote{авторизация на точную сумму, клиринг на ту же сумму,
расчёт через день} описывает идеальный случай и покрывает основную
массу розничных покупок. Платёжная индустрия полна сценариев,
которые в~эту модель не укладываются, и~для каждого есть отдельный
механизм.

\subsection*{Предавторизация (pre-authorization)}

Предавторизация, \textit{pre-auth}, -- авторизация, при которой финальная
сумма заранее неизвестна. Классические примеры -- гостиница и
аренда автомобиля. В~момент заселения или получения машины продавец
авторизует оценочную сумму -- ожидаемую стоимость проживания
или депозит за автомобиль, -- а~итоговая определяется при выезде
или возврате и~может оказаться как меньше, так и~больше
авторизованной.

В~сетевых правилах \textit{pre-auth} -- самостоятельный тип запроса. Visa
различает Estimated Authorization и Incremental Authorization,
а Mastercard рекомендует кодировать запрос как \textit{preauthorization},
если сумма оценочная или если операция может вообще не завершиться.
Эмитент ставит холд, заранее зная, что финальный клиринг может
отличаться от исходной оценки. Допустимые пределы расхождения
зависят от MCC, типа ТСП и конкретного сценария, поэтому точные
пороги берут из актуальных правил сети~\cite{visa-tadg,visa-core-rules,mc-tpr}.

\practicenote{%
  Если вы проектируете систему бронирования или маркетплейс услуг
  с~переменной стоимостью, \textit{pre-auth} становится вашим основным инструментом.
  Но учтите два нюанса. Во-первых, холд на карте клиента будет
  стоять на полную авторизованную сумму, даже если итоговая покажется
  меньше. Во-вторых, если финальная сумма окажется значительно
  больше авторизованной, вам потребуется либо инкрементальная
  авторизация, либо новая транзакция на разницу.

}
\subsection*{Инкрементальная авторизация (incremental authorization)}

Инкрементальная авторизация увеличивает сумму ранее одобренной
авторизации без отмены исходной. Продавец отправляет дополнительный
авторизационный запрос, ссылающийся на исходный через код авторизации
или специальный идентификатор, и~эмитент увеличивает холд
на запрошенную сумму~\cite{visa-tadg,visa-core-rules,mc-tpr}.

Типичный сценарий -- гостиница, где гость продлевает пребывание
или заказывает дополнительные услуги. Вместо повторного цикла
\enquote{авторизовать -- отменить} продавец наращивает сумму
инкрементально. Поддержка такого потока зависит от схемы, региона,
MCC и~настроек эмитента или эквайера, поэтому при проектировании
интеграции всегда нужен запасной сценарий -- скажем, новая
авторизация на дополнительную сумму.

\subsection*{Частичная авторизация (partial authorization)}

Частичная авторизация -- сценарий, в~котором эмитент одобряет
транзакцию не на полную запрошенную сумму, а~лишь на часть. Держатель
пытается оплатить покупку на тысячу рублей, но на карте доступно
только семьсот. Вместо полного отказа с~кодом \texttt{51} эмитент
одобряет операцию на семьсот рублей, а~продавец предлагает клиенту
доплатить оставшиеся триста наличными или другой
картой~\cite{visa-tadg,mc-tpr}.

Частичная авторизация требует поддержки на всех уровнях. Терминал
обрабатывает частичное одобрение и предлагает доплату, POS-система
корректирует чек, кассир должен понимать, что происходит. Чаще всего
схема встречается с~предоплаченными и подарочными картами, где остаток
заранее известен и ограничен.

\importantnote{%
  Если ваша система не поддерживает partial authorization, убедитесь,
  что терминал или платёжный шлюз не заявляет такую поддержку
  в~сетевом индикаторе partial approval. Например, у
  Mastercard для этого используется Partial Approval
  Terminal Support Indicator в DE~48, subelement~61.
  Получить частичное одобрение и не обработать его корректно --
  верный путь к~расхождениям в~учёте, когда продавец отпустит товар
  на полную сумму, а~спишется только часть~\cite{mc-tpr}.

}
\subsection*{Отложенная авторизация (deferred authorization)}

Отложенная авторизация -- механизм, при котором транзакция совершается
без онлайн-авторизации, а авторизационный запрос отправляется позже,
по восстановлении связи. Классические сценарии -- оплата на борту
самолёта, где терминал работает офлайн и авторизации уходят после
посадки; паркинг, где сумма определяется по времени стоянки;
платные дороги и общественный транспорт.

В~режиме deferred authorization терминал собирает данные карты,
фиксирует операцию локально и~позже формирует авторизационный запрос
с~соответствующим флагом. Эмитент может одобрить или отклонить
этот запрос, но товар или услуга уже предоставлены, и~кредитный риск
лежит на эквайере и~продавце. Поэтому deferred authorization разрешена
только для определённых MCC и~при соблюдении лимитов на сумму
отдельной операции, устанавливаемых сетями. У Visa такие авторизации
нужно завершить в~течение 24~часов с~момента исходной операции,
а для авиаперевозок и ряда транспортных MCC -- 4111, 4112, 4131 --
допускается до 4 суток. В~авторизационном сообщении передаётся
Deferred Authorization Indicator со значением \texttt{5206} в~поле
Field~63.3, обязательный и для card-present, и~для card-absent
сценариев~\cite{visa-tadg,visa-core-rules}.

\subsection*{Принудительная отправка (force post)}

Принудительная отправка -- механизм, при котором продавец отправляет
транзакцию в~клиринг без предварительной онлайн-авторизации, используя
код авторизации, полученный другим путём -- как правило, по телефону
через голосовую авторизацию, voice authorization.

Сценарий типичен, когда онлайн-авторизация невозможна -- терминал
не работает, связь пропала, -- но провести платёж нужно. Продавец
получает голосовую авторизацию у~эмитента или через эквайера
и~затем использует этот код при обработке транзакции. Mastercard
прямо оговаривает -- код, полученный по каналу голосовой авторизации,
сам по себе не защищает от чарджбэка по авторизационным
основаниям~\cite{mc-tpr}.

Force post несёт повышенные риски. Онлайн-проверки криптограммы
здесь нет, часть контроля переносится на человека и процедуры
эквайера, и ответственность за последствия в~значительной степени
ложится на продавца. Сам механизм становится всё более редким --
устойчивые каналы связи и~офлайн-режимы терминалов покрывают
большинство сценариев, где раньше требовалась голосовая авторизация.
Полностью force post не исчезает, но живёт как запасной вариант.

\begin{table}[ht]
  \begingroup
  \small
  \setlength{\tabcolsep}{4pt}
  \renewcommand{\arraystretch}{1.2}
  \makebox[\textwidth][l]{%
    \begin{tabularx}{\fullwidth}{
        @{}>{\RaggedRight\arraybackslash\bfseries}p{0.14\textwidth}
        >{\RaggedRight\arraybackslash}p{0.17\textwidth}
        >{\RaggedRight\arraybackslash}p{0.15\textwidth}
        >{\RaggedRight\arraybackslash}p{0.17\textwidth}
        Y@{}
      }
      \toprule
      \rowcolor{Panel}
      \textbf{Механизм} & \textbf{Когда сумма известна} &
      \textbf{Онлайн-авторизация} & \textbf{Кто несёт риск} &
      \textbf{Типичный MCC} \\
      \midrule
      \rowcolor{Soft}
      Pre-auth & позже & да
      & эквайер/ТСП & 7011 (гостиницы), 7512 (аренда авто) \\
      Incremental auth & постепенно & да (несколько)
      & эквайер/ТСП & 7011 \\
      \rowcolor{Soft}
      Partial auth & сразу, но покрытие неполное & да
      & ТСП & любой (чаще prepaid) \\
      Deferred auth & позже & нет (офлайн)
      & эквайер/ТСП & 4511 (авиа), 7523 (паркинг) \\
      \rowcolor{Soft}
      Force post & сразу & голосовая
      & ТСП & любой \\
      \bottomrule
    \end{tabularx}%
  }
  \endgroup
  \caption{Граничные случаи карточной авторизации.}\label{fig:tx-edge-cases}
\end{table}
